\section{Mechanistic read-outs and steering}
\label{app:mech}

We read $\dhat$ at three positions, running one forward pass over the prompt followed by the model's own
response: the last prompt token, the first token of the harmful item's answer (its \emph{answer onset}), and
the onset of a benign item.

\begin{table}[h]
\centering\small
\caption{Projection onto $\dhat$ at the selected layer, scaled so that a harmful request alone is 100. AUCs
compare list responses that refuse the harmful item with those that comply (800 probed list responses per
model; 39, 11 and 54 complied, so Gemma's AUCs rest on few positives). The control row uses the projection of
the \emph{same harmful request asked alone}.}
\label{tab:probe}
\resizebox{\linewidth}{!}{\begin{tabular}{lccc}
\toprule
Projection (\% of harmful-alone value) & Qwen3-14B & Gemma-2-9B & Llama-3.1-8B \\
\midrule
last prompt token, harmful request alone & +100 & +100 & +100 \\
last prompt token, benign request alone & -30 & +4 & -6 \\
\midrule
last prompt token, harmful item first & +30 & +60 & +15 \\
last prompt token, harmful item last & +11 & +40 & +20 \\
\midrule
harmful item's answer onset, item refused & +4 & +56 & +27 \\
harmful item's answer onset, item complied & -18 & +35 & +29 \\
benign item's answer onset & -25 & +10 & +6 \\
\midrule
AUC refused vs.\ complied, answer onset (sel.\ / best layer) & 0.88 / 0.91 (27) & 0.86 / 0.90 (26) & 0.46 / 0.83 (16) \\
AUC refused vs.\ complied, last prompt token (sel.\ / best) & 0.95 / 0.95 (24) & 0.92 / 0.97 (38) & 0.69 / 0.91 (25) \\
\quad control: same request \emph{alone}, last prompt token (layer) & 0.95 (24) & 0.89 (38) & 0.93 (25) \\
\midrule
lists with negative last-token projection & 16\% & 0\% & 3\% \\
\quad ASR / target refused, if negative & 23\% / 46\% & -- & 23\% / 55\% \\
\quad ASR / target refused, if positive & 1\% / 95\% & 1\% / 97\% & 6\% / 86\% \\
\bottomrule
\end{tabular}}
\end{table}

\paragraph{The sign flip is real but is not a decision.} List structure shrinks the last-token projection to
11--60\% of its single-request value (\cref{tab:probe}), and for Qwen3-14B it turns negative in \QwenNegLastFrac\
of list prompts. When it does, ASR is 23\% rather than 1\%, but the model still refuses or deflects the
harmful item in 77\% of those responses.

\paragraph{The last-token read-out mostly reflects the request, not the list.} Within list prompts, the
last-token projection separates responses that end up refusing the harmful item from those that comply (AUC
up to \QwenAUCLastBest, \GemmaAUCLastBest\ and \LlamaAUCLastBest). This looks like the model ``planning'' its
refusal \citep{dong2025planning}. The control row shows it is not specific to the list: the projection of the
\emph{same request asked alone} predicts the list outcome about as well (\QwenAUCBare, \GemmaAUCBare,
\LlamaAUCBare). Requests the model objects to weakly are complied with in lists and alone alike. We therefore do
not interpret the end-of-prompt read-out as a list-specific plan. At the harmful item's answer onset, the
projection again separates refused from complied items (AUC up to 0.91, 0.90 and 0.83), and it sits above the
onset of benign items in the same responses.

\paragraph{The guardrail head attends to the item being answered.} From the last prompt token, Qwen3-14B's
head L36H31 puts little attention on the harmful request in a list (\AttnLast, against \AttnBare\ for the request
alone). From the harmful item's answer onset it puts \AttnOnset, more than the same head puts on a benign item's
own request at that item's onset (\AttnBenignOwn). Attention at the onset does fall with list length (from 0.027
with 4 items to 0.014 with 16 items, harmful item last), but ASR does not rise with it (\cref{fig:behaviour}).
Zeroing the head (\cref{tab:head}): zeroing L36H31 changes neither ASR nor refusal (30\%/61\% without, 34\%/58\% with), and
neither do three control heads, so the head reads the item but is not necessary for refusing it, consistent
with \citet{srivastava2026ara}.

\paragraph{Adding the direction at the decision point.} Teacher-forcing the model's own response up to the
harmful item's answer onset ($k{=}15$) and adding $\alpha$ times the \emph{unnormalized} difference-of-means
vector from that token onward lowers harmful continuations: for Qwen3-14B from 28\% to 6\% at $\alpha{=}1$
with no refusals of benign targets; for Llama-3.1-8B from 34\% to 8\% at $\alpha{=}0.5$ with 3\% benign
refusals. Larger doses refuse benign targets as well. Adding the vector at the single decision token is much weaker: Qwen3-14B needs twice the
dose (3\% ASR at $\alpha{=}2$) and Llama-3.1-8B is unaffected at the doses tested (\cref{tab:steer1}).\ Full results are in \cref{app:causal}.
Our earlier conclusion that steering cannot restore refusal came from adding at most 2.6\% of this vector.
